% !TEX root = task-6-1-3.tex
% Editable task bank. Append new BankTask blocks at the end of the document.
\documentclass[12pt,a4paper]{article}
\usepackage{physics-textbook}
\usepackage{graphicx,needspace}
\usetikzlibrary{decorations.pathmorphing}
\newif\ifShowSolutions
\ifdefined\TutorEdition
  \ShowSolutionstrue
\else
  \ShowSolutionsfalse % Change to \ShowSolutionstrue to display solutions.
\fi
\newcounter{banktask}
\newsavebox{\BankTaskBox}
\newcommand{\BankTaskContent}[3]{%
  \noindent\textbf{\thebanktask.}\quad #2\par
  \ifShowSolutions
    \Needspace{50mm}% Keep the solution heading with its opening text or diagram.
    #3\par
    \ifdefined\TutorEdition
      \ifcsname TutorSource#1\endcsname
        \nopagebreak\smallskip{\small\csname TutorSource#1\endcsname\par}%
      \fi
    \fi
  \fi
}
\newcommand{\BankTask}[3]{%
  \refstepcounter{banktask}\par\medskip
  \sbox{\BankTaskBox}{\begin{minipage}{\linewidth}%
    \BankTaskContent{#1}{#2}{#3}\end{minipage}}%
  \ifdim\dimexpr\ht\BankTaskBox+\dp\BankTaskBox\relax>\dimexpr\textheight-20mm\relax
    \Needspace{35mm}\BankTaskContent{#1}{#2}{#3}%
  \else
    \Needspace{\dimexpr\ht\BankTaskBox+\dp\BankTaskBox\relax}%
    \noindent\usebox{\BankTaskBox}\par
  \fi
  \medskip
}
\begin{document}
\BookSubsection{6.1.3}{Законы фотоэффекта}


% Source sdamgia-3116: https://phys-ege.sdamgia.ru/problem?id=3116
\BankTask{sdamgia-3116}{
Металлическую пластину освещали монохроматическим светом с длиной волны $\lambda =500$ нм. Что произойдет с частотой падающего света, импульсом фотонов и кинетической энергией вылетающих электронов при освещении этой пластины монохроматическим светом с длиной волны $\lambda =700$ нм одинаковой интенсивности? Фотоэффект наблюдается в обоих случаях.
$1$)  увеличилась;
$2$)  уменьшилась;
$3$)  не изменилась.

\begin{center}\begin{tabular}{|p{.29\linewidth}|p{.29\linewidth}|p{.29\linewidth}|}\hline Частота падающего света & Импульс фотонов & Кинетическая энергия фотоэлектронов\\\hline\strut&\strut&\strut\\\hline\end{tabular}\end{center}
}{
\smallskip\noindent\textbf{Решение}\par
Частота света связана с длиной волны и скоростью света соотношением $\lambda \nu =c.$ Следовательно, увеличение длины волны падающего света соответствует уменьшению частоты (A  — $2$). Импульс фотона обратно пропорционален длине его волны: $p= \frac{h}{\lambda} .$ Таким образом, при увеличении длины волны, импульс фотонов уменьшается (Б  — $2$). Кинетическая энергия вылетающих электронов связана с энергией фотонов и работой выхода, согласно уравнению фотоэффекта, соотношением 

\[E_{\,\text{фот}}= \frac{hc}{\lambda} =A_{\,\text{вых}} + E_{\,\text{кин}}.\]

Работа выхода зависит только от химических свойств металлов, а значит, в результате увеличения длины кинетическая энергия фотоэлектронов уменьшится (В  — $2$).

Ответ: $222$.
}

% Source sdamgia-3158: https://phys-ege.sdamgia.ru/problem?id=3158
\BankTask{sdamgia-3158}{
Квант света выбивает электрон из металла. Как изменятся при увеличении энергии фотона в этом опыте следующие три величины: работа выхода электрона из металла, максимальная возможная скорость фотоэлектрона, его максимальная кинетическая энергия? 

$1$)  увеличится;
$2$)  уменьшится;
$3$)  не изменится.

\begin{center}\begin{tabular}{|p{.29\linewidth}|p{.29\linewidth}|p{.29\linewidth}|}\hline Работа выхода электрона из металла & Максимальная скорость фотоэлектрона & Максимальная кинетическая энергия\\\hline\strut&\strut&\strut\\\hline\end{tabular}\end{center}
}{
\smallskip\noindent\textbf{Решение}\par
Работа выхода электрона из металла зависит только от свойств металла и никак не изменяется при увеличении энергии фотона в этом опыте. Согласно уравнению фотоэффекта Эйнштейна, энергия поглощенного фотона идет на работу выхода электрона из металла и на сообщение этому электрону кинетической энергии: $E_{\,\text{фот}}=A_{\,\text{вых}} + E_{\,\text{кин}}.$ Следовательно, при увеличении энергии фотона максимальная кинетическая энергия фотоэлектронов увеличивается. Раз увеличивается максимальная кинетическая энергия, увеличивается и максимальная скорость фотоэлектронов.

Ответ: $311$.
}

% Source sdamgia-9331: https://phys-ege.sdamgia.ru/problem?id=9331
\BankTask{sdamgia-9331}{
\begin{center}\begin{tikzpicture}[book diagram,>=stealth]\draw[task axis,->](-1.4,0)--(3.8,0)node[right]{$U,\,\text{В}$};\draw[task axis,->](0,-.2)--(0,2.5)node[above]{$I,\,\text{мА}$};\node[below left]at(0,0){$0$};\draw[task curve]( -1,0)..controls(0,0)and(.15,.35)..(.6,1.2)..controls(1.05,2.05)and(1.65,2)..(2,2)--(3.2,2);\draw[task projection](0,2)--(2,2);\foreach\x in{-1,1,2,3}{\draw(\x,.05)--(\x,-.05);\node[below]at(\x,0){$\x$};}\foreach\y in{1,2}{\draw(-.05,\y)--(.05,\y);\node[left]at(0,\y){$\y$};}\end{tikzpicture}\end{center}
В опыте по изучению фотоэффекта свет частотой $\nu$ = $6{,}1$ · $10^{14}$ Гц падает на поверхность катода, в результате чего в цепи возникает ток. График зависимости силы тока I от напряжения U между анодом и катодом приведен на рисунке. Какова мощность падающего света P, если в среднем один из $20$ фотонов, падающих на катод, выбивает электрон?
}{
\smallskip\noindent\textbf{Решение}\par
Из графика находим величину тока насыщения, которая равна $2$ мА. Ток насыщения соответствует максимальному потоку электронов, которое способно выбивать в единицу времени излучение мощностью $P.$
По определению, сила тока  — это количество заряда, прошедшего за единицу времени:

\[I= \frac{q}{t} = \frac{N_e |e|}{t} .\]

Мощность светового потока  — это энергия, которую несут фотоны за единицу времени:

\[P= \frac{W}{t} = \frac{N_{\,\text{ф}} h\nu}{t} .\]

Учтем, что один электрон выбивается каждые $20$ фотонов, то есть $N_{\,\text{ф}}=20N_e:$

\[P= \frac{20N_e h\nu}{t} = \frac{20I_{\,\text{н}}h\nu}{|e|} = \frac{20 \cdot 2 \cdot 10 ^{ - 3 } \cdot 6{,}6 \cdot 10 ^{ - 34 } \cdot 6{,}1 \cdot 10 ^{ 14 }}{1{,}6 \cdot 10 ^{ - 19 }} \approx0{,}1\,\text{Вт}.\]

Ответ: $0{,}1$ Вт
}

% Source sdamgia-27989: https://phys-ege.sdamgia.ru/problem?id=27989
\BankTask{sdamgia-27989}{
Для проведения опытов по наблюдению фотоэффекта взяли пластину из металла с работой выхода $3{,}4\cdot10^{-19}\,\text{Дж}$ и стали освещать её светом частотой $6\cdot10^{14}\,\text{Гц}$. Как изменится модуль запирающего напряжения $U_{\text{з}}$ и максимальный импульс фотоэлектронов $p_{\max}$, вылетающих с поверхности металла, если увеличить интенсивность падающего света, не изменяя его частоту?
\par1. Увеличится. 2. Уменьшится. 3. Не изменится.
\TaskChangeTable{Модуль запирающего напряжения $U_{\text{з}}$}{Максимальный импульс фотоэлектронов $p_{\max}$}
}{
\smallskip\noindent\textbf{Решение}\par
При увеличении интенсивности падающего света увеличится количество фотонов, но при неизменности частоты падающего света максимальная кинетическая энергия вылетающих электронов не изменится. Следовательно, не изменится модуль запирающего напряжения ($3$).
Поскольку не изменится максимальная кинетическая энергия фотоэлектронов, то не изменится и их импульс ($3$).

Ответ: $33$.
}

% Source sdamgia-7193: https://phys-ege.sdamgia.ru/problem?id=7193
\BankTask{sdamgia-7193}{
На металлическую пластинку направили пучок света от лазера, вызвав фотоэффект. Интенсивность лазерного излучения плавно увеличивают, не меняя его частоты. Как меняются в результате этого число вылетающих в единицу времени фотоэлектронов и их максимальная кинетическая энергия?

$1$.  Увеличится.
$2$.  Уменьшится.
$3$.  Не изменится.

 

\TaskChangeTable{Число фотоэлектронов, вылетающих в единицу времени}{Максимальная кинетическая энергия фотоэлектронов}
}{
\smallskip\noindent\textbf{Решение}\par
Интенсивность лазерного излучения  — количество фотонов, прошедших через единицу площади за единицу времени. Значит, при увеличении интенсивности число фотонов увеличится и увеличится число фотоэлектронов. По второму закону фотоэффекта $E_{\,\text{фот}}=A_{\,\text{вых}} + E_{\,\text{кин}}$ кинетическая энергия фотоэлектронов зависит от работы выхода материала и от энергии фотона. При увеличении интенсивности энергия фотонов не изменяется, а значит, максимальная кинетическая энергия фотоэлектронов остается неизменной.

Ответ: $13$.
}

% Source sdamgia-11563: https://phys-ege.sdamgia.ru/problem?id=11563
\BankTask{sdamgia-11563}{
В первом опыте по изучению фотоэффекта металлическую пластинку освещают белым светом через синий светофильтр (пропускает только синий цвет), а во втором  — через зеленый (пропускает только зеленый цвет). Как изменяются следующие величины при переходе от первого опыта ко второму?

$1$.  Увеличилась.
$2$.  Уменьшилась.
$3$.  Не изменилась.

\TaskChangeTable{Частота падающего на пластинку света}{Работа выхода электронов из металла}
}{
\smallskip\noindent\textbf{Решение}\par
Из формулы $h\nu=A_{\,\text{вых}} + E_{\text{кин,max}}$ можно сделать вывод, что частота падающего на пластинку света уменьшилась, а работа выхода электронов из металла не изменилась.

Ответ: $23$.
}

% Source sdamgia-6159: https://phys-ege.sdamgia.ru/problem?id=6159
\BankTask{sdamgia-6159}{
\begin{center}\begin{tikzpicture}[book diagram,>=stealth]\draw[task axis,->](-.2,0)--(3.8,0)node[right]{$U,\,\text{В}$};\draw[task axis,->](0,-.2)--(0,2.5)node[above]{$I,\,\text{мА}$};\node[below left]at(0,0){$0$};\draw[task curve](0,2.2)..controls(.2,1.1)and(.9,.25)..(2,0);\foreach\x in{1,2}{\fill(\x,0)circle(1.4pt);\node[below]at(\x,0){$\x$};}\end{tikzpicture}\end{center}
Работа выхода для некоторого металла равна $3\,\text{эВ}$. На пластинку из этого металла падает свет. На рисунке показана зависимость силы I фототока от приложенного обратного напряжения U. Какова энергия фотона светового излучения, падающего на эту пластинку? (в электрон-вольтах.)
}{
\smallskip\noindent\textbf{Решение}\par
При вылете фотоэлектрона часть энергии фотона затрачивается на преодоление работы выхода металла, а оставшаяся часть уходит на увеличение кинетической энергии электрона $h\nu=A_{\,\text{вых}} + E_{\,\text{кин}}.$ При увеличении запирающего напряжения фототок будет уменьшаться, в связи с тем, что кинетической энергии электронов не будет хватать для преодоления электрического поля. Фототок исчезнет, когда потенциальная энергия электронов в электрическом поле станет равной их кинетической энергии $E_{\,\text{кин}}=Ue.$ Из рисунка видим, что фототок исчезает при напряжении $2\,\text{В}$. Следовательно, максимальная энергия фотоэлектронов $2\,\text{эВ}$. Значит, энергия фотона светового излучения равна $2$ + $3$  =  $5\,\text{эВ}$.

Ответ: $5$.
}

% Source sdamgia-5488: https://phys-ege.sdamgia.ru/problem?id=5488
\BankTask{sdamgia-5488}{
\begin{center}\begin{tikzpicture}[book diagram,>=stealth]\fill(-2,0)circle(.06)node[above left]{$S$};\draw[<->,thick](0,-1.3)--(0,1.3)node[above]{Линза};\draw[thick](2,-1.3)--(.7,1.3)node[above right]{Фотокатод};\foreach\y in{-1,-.5,0,.5,1}{\draw[->](-2,0)--(0,\y);\draw[->](0,\y)--({1.35-.5*\y},\y);}\end{tikzpicture}\end{center}
В установке по наблюдению фотоэффекта свет от точечного источника S, пройдя через собирающую линзу, падает на фотокатод параллельным пучком. В схему внесли изменение: на место первоначальной линзы поставили собирающую линзу того же диаметра, но с меньшим фокусным расстоянием. Источник света переместили вдоль главной оптической оси линзы так, что на фотокатод свет снова стал падать параллельным пучком. Как изменился при этом (уменьшился или увеличился) фототок насыщения? Объясните, почему изменяется фототок насыщения, и укажите, какие физические закономерности вы использовали для объяснения.
}{
\smallskip\noindent\textbf{Решение}\par
По первому закону Столетова фототок насыщения зависит от интенсивности падающего света, то есть от количества фотонов, падающих на фотокатод в единицу времени. При использовании линзы такого же диаметра, но с меньшим фокусным расстоянием, телесный угол, под которым из источника видно линзу, увеличивается, поскольку источник теперь расположен ближе к ней (для получения параллельного пучка источник нужно разместить в фокусе линзы). Фотоны летят от источника во все стороны равномерно, поэтому результирующий поток фотонов, попадающих на фотокатод в результате замены линзы, увеличивается. А значит, увеличивается и ток насыщения.
}

% Source sdamgia-7933: https://phys-ege.sdamgia.ru/problem?id=7933
\BankTask{sdamgia-7933}{
В опыте по изучению фотоэффекта катод освещается желтым светом, в результате чего в цепи возникает ток (см. рис. $1$). Зависимость показаний амперметра I от напряжения U между анодом и катодом приведена на втором рисунке. Используя законы фотоэффекта и предполагая, что отношение числа фотоэлектронов к числу поглощенных фотонов не зависит от частоты света, объясните, как изменится представленная зависимость I(U), если освещать катод зеленым светом, оставив мощность поглощенного катодом света неизменной. 

\begin{center}\begin{tikzpicture}[book diagram,>=stealth]\draw(0,1.4)ellipse(1.4 and .8);\draw(-1.1,.9)--(-1.1,1.9) (1.1,.9)--(1.1,1.9);\draw(-1.1,1.4)--(-2,1.4)--(-2,-1)--(-1.2,-1);\draw(-.9,-1)circle(.3)node{$A$};\draw(-.6,-1)--(.8,-1) (1,-1)--(2,-1)--(2,1.4)--(1.1,1.4);\draw[thick](.8,-1.15)--(.8,-.85) (1,-1.3)--(1,-.7);\node[below]at(.9,-1.35){$\mathcal E$};\foreach\y in{1.2,1.55,1.9}{\draw[->]({-.3+2*(\y-1.2)},2.5)--(-1.1,\y);}\node[above]at(.4,2.5){свет};\end{tikzpicture}\end{center}

Рис. $1$

\begin{center}\begin{tikzpicture}[book diagram,>=stealth]\draw[task axis,->](-1.4,0)--(3.8,0)node[right]{$U$};\draw[task axis,->](0,-.2)--(0,2.5)node[above]{$I$};\node[below left]at(0,0){$0$};\draw[task curve]( -1,0)..controls(0,0)and(.15,.35)..(.6,1.2)..controls(1.05,2.05)and(1.65,2)..(2,2)--(3.2,2);\draw[task projection](0,2)--(2,2);\node[below]at(-1,0){$U_{\text{з}}$};\node[left]at(0,2){$I_{\text{нас}}$};\end{tikzpicture}\end{center}

Рис. $2$
}{
\smallskip\noindent\textbf{Решение}\par
$1$.  При изменении света с желтого на зеленый его длина волны уменьшится, частота увеличится ($\nu_{\,\text{з}}$ > $\nu_{\,\text{ж}}$). 
$2$.  Работа выхода электронов из материала не зависит от частоты падающего света, поэтому в соответствии с уравнением Эйнштейна для фотоэффекта: $h\nu$ = $A_{\,\text{вых}}$ + $E_{max}$  — увеличится максимальная кинетическая энергия фотоэлектронов $E_{max}$. Так как $E_{\max}=e|U_{\,\text{з}}|,$ то увеличится и модуль запирающего напряжения $U_{\,\text{з}}$. 
$3$.  Мощность поглощенного света связана с частотой волны $\nu$ соотношением P  =  $N_{\varphi }$$E_{\varphi }$  =  $N_{\varphi }$h$\nu$, где $N_{\varphi }$  — число фотонов, падающих на катод за $1\,\text{с}$, $E_{\varphi }$= h$\nu$ — энергия одного фотона (соотношение Планка). Так как мощность света не изменилась, а энергия фотонов $E_{\varphi }$ увеличилась, то уменьшится число фотонов, падающих на катод за $1\,\text{с}$.
$4$.  Сила тока насыщения $I_{\,\text{нас}}$ определяется числом выбитых светом за $1\,\text{с}$ электронов $N_{e}$, которое пропорционально числу падающих на катод за $1\,\text{с}$ фотонов, поэтому сила тока насыщения уменьшится. 

Ответ: точка отрыва графика от горизонтальной оси U сдвинется влево, горизонтальная асимптота графика $I_{\,\text{нас}}$ сдвинется вниз.
}

% Source sdamgia-32265: https://phys-ege.sdamgia.ru/problem?id=32265
\BankTask{sdamgia-32265}{
В классе при электрическом освещении лампами накаливания показали опыт: цинковый шар электрометра зарядили эбонитовой палочкой, потертой о сукно. При этом стрелка электрометра отклонилась, заняв положение, указанное на рисунке, и не меняла его. Через некоторое время к шару на расстояние нескольких сантиметров поднесли горящую свечу, при этом 
 стрелка электрометра быстро опустилась вниз. Спектр излучения свечи (зависимость интенсивности излучения I от длины волны) показан на рисунке. Объясните разрядку электрометра, примите во внимание приведенные спектры и то, что для цинка «красная граница» фотоэффекта $\lambda_{\,\text{кр}}$  =  $290\,\text{нм}$.

\begin{center}\begin{tikzpicture}[book diagram,>=stealth]\draw(0,0)circle(1);\draw(0,0)circle(.91);\draw(0,1)--(0,1.6);\draw[fill=white](0,2.05)circle(.5);\draw(0,.9)--(0,-.85);\draw[very thick](-.75,-.5)--(.7,.48);\foreach\a in{220,230,...,270}{\draw(\a:.75)--(\a:.9);}\draw(-.25,-1)--(-.4,-1.5)--(.4,-1.5)--(.25,-1);\draw(.7,-.7)--(1.7,-.7)--(1.7,-1.2);\draw(1.5,-1.2)--(1.9,-1.2) (1.57,-1.3)--(1.83,-1.3) (1.64,-1.4)--(1.76,-1.4);\node[right]at(1.8,-1){Заземление};\end{tikzpicture}\end{center}

\begin{center}\begin{tikzpicture}[book diagram,>=stealth]\begin{scope}[x=.010cm,y=.15cm]\draw[task axis,->](350,0)--(1100,0);\draw[task axis,->](350,0)--(350,27);\draw[task curve]plot coordinates{(350,0)(353.28,0.20000)(356.55,0.23300)(359.83,0.30000)(363.11,0.40000)(366.39,0.50000)(369.66,0.56700)(372.94,0.70000)(376.22,0.76700)(379.5,0.90000)(382.77,0.96700)(386.05,1.10000)(389.33,1.23300)(392.61,1.30000)(395.88,1.43300)(399.16,1.56700)(402.44,1.70000)(405.71,1.83300)(408.99,1.96700)(412.27,2.10000)(415.55,2.26700)(418.82,2.43300)(422.1,2.60000)(425.38,2.76700)(428.66,2.93300)(431.93,3.13300)(435.21,3.33300)(438.49,3.53300)(441.76,3.80000)(445.04,4.03300)(448.32,4.30000)(451.6,4.56700)(454.87,4.90000)(458.15,5.23300)(461.43,5.60000)(464.71,6.03300)(467.98,6.50000)(471.26,7.03300)(474.54,7.60000)(477.82,8.23300)(481.09,8.93300)(484.37,9.70000)(487.65,10.60000)(490.92,11.73300)(494.2,12.83300)(497.48,13.80000)(500.76,14.43300)(504.03,14.90000)(507.31,15.30000)(510.59,15.63300)(513.87,15.90000)(517.14,16.13300)(520.42,16.33300)(523.7,16.53300)(526.97,16.70000)(530.25,16.86700)(533.53,17.03300)(536.81,17.16700)(540.08,17.23300)(543.36,17.36700)(546.64,18.76700)(549.92,19.33300)(553.19,19.00000)(556.47,18.76700)(559.75,18.63300)(563.03,18.56700)(566.3,18.53300)(569.58,18.60000)(572.86,18.70000)(576.13,18.83300)(579.41,19.06700)(582.69,19.26700)(585.97,19.53300)(589.24,19.76700)(592.52,20.06700)(595.8,20.36700)(599.08,20.73300)(602.35,20.93300)(605.63,20.56700)(608.91,20.16700)(612.18,19.76700)(615.46,19.36700)(618.74,18.96700)(622.02,18.60000)(625.29,18.26700)(628.57,17.86700)(631.85,17.53300)(635.13,17.20000)(638.4,16.80000)(641.68,16.50000)(644.96,16.16700)(648.24,15.83300)(651.51,15.53300)(654.79,15.20000)(658.07,14.86700)(661.34,14.56700)(664.62,14.30000)(667.9,13.96700)(671.18,13.70000)(674.45,13.36700)(677.73,13.10000)(681.01,12.83300)(684.29,12.53300)(687.56,12.26700)(690.84,12.00000)(694.12,11.73300)(697.39,11.50000)(700.67,11.26700)(703.95,11.00000)(707.23,10.76700)(710.5,10.53300)(713.78,10.33300)(717.06,10.06700)(720.34,9.86700)(723.61,9.66700)(726.89,9.40000)(730.17,9.20000)(733.45,9.00000)(736.72,8.80000)(740.0,8.60000)(743.28,8.40000)(746.55,8.20000)(749.83,8.00000)(753.11,7.83300)(756.39,7.66700)(759.66,7.46700)(762.94,7.30000)(766.22,7.13300)(769.5,6.96700)(772.77,6.80000)(776.05,6.63300)(779.33,6.50000)(782.61,6.36700)(785.88,6.23300)(789.16,6.10000)(792.44,5.96700)(795.71,5.83300)(798.99,5.76700)(802.27,5.63300)(805.55,5.53300)(808.82,5.43300)(812.1,5.36700)(815.38,5.30000)(818.66,5.20000)(821.93,5.13300)(825.21,5.06700)(828.49,5.00000)(831.76,4.93300)(835.04,4.90000)(838.32,4.83300)(841.6,4.80000)(844.87,4.76700)(848.15,4.73300)(851.43,4.70000)(854.71,4.66700)(857.98,4.63300)(861.26,4.60000)(864.54,4.56700)(867.82,4.53300)(871.09,4.50000)(874.37,4.46700)(877.65,4.40000)(880.92,4.33300)(884.2,4.30000)(887.48,4.23300)(890.76,4.16700)(894.03,4.10000)(897.31,4.03300)(900.59,3.96700)(903.87,3.90000)(907.14,3.83300)(910.42,3.73300)(913.7,3.66700)(916.97,3.56700)(920.25,3.50000)(923.53,3.36700)(926.81,3.30000)(930.08,3.20000)(933.36,3.10000)(936.64,3.00000)(939.92,2.90000)(943.19,2.76700)(946.47,2.70000)(949.75,2.56700)(953.03,2.46700)(956.3,2.36700)(959.58,2.23300)(962.86,2.10000)(966.13,2.00000)(969.41,1.90000)(972.69,1.76700)(975.97,1.63300)(979.24,1.50000)(982.52,1.36700)(985.8,1.23300)(989.08,1.16700)};\node[below]at(400,0){$ 400 $};\node[below]at(500,0){$ 500 $};\node[below]at(600,0){$ 600 $};\node[below]at(700,0){$ 700 $};\node[below]at(800,0){$ 800 $};\node[below]at(900,0){$ 900 $};\node[below]at(1000,0){$ 1000 $};\node[left]at(360,10.0){$ 100 $};\node[left]at(360,20.0){$ 200 $};\node[above]at(360,27){$I,\,\text{отн. ед.}$};\node[above,align=center,text width=5cm]at(760,23){Спектр излучения лампы накаливания};\node[below]at(750,-4){Длина волны, нм};\end{scope}\end{tikzpicture}\end{center}

\begin{center}\begin{tikzpicture}[book diagram,>=stealth]\begin{scope}[x=.010cm,y=.15cm]\draw[task axis,->](350,0)--(1100,0);\draw[task axis,->](350,0)--(350,27);\draw[task curve]plot coordinates{(350,0)(353.28,0.11750)(356.55,0.14750)(359.83,0.14750)(363.11,0.17750)(366.39,0.17750)(369.66,0.23500)(372.94,0.26500)(376.22,0.29500)(379.5,0.32500)(382.77,0.35500)(386.05,0.41250)(389.33,0.44250)(392.61,0.47250)(395.88,0.53000)(399.16,0.56000)(402.44,0.59000)(405.71,0.65000)(408.99,0.70750)(412.27,0.73750)(415.55,0.79750)(418.82,0.82500)(422.1,0.88500)(425.38,0.94500)(428.66,1.00250)(431.93,1.06250)(435.21,1.12000)(438.49,1.18000)(441.76,1.24000)(445.04,1.29750)(448.32,1.35750)(451.6,1.41500)(454.87,1.50500)(458.15,1.56250)(461.43,1.62250)(464.71,1.71000)(467.98,1.80000)(471.26,1.85750)(474.54,1.94750)(477.82,2.03500)(481.09,2.15250)(484.37,2.21250)(487.65,2.33000)(490.92,2.44750)(494.2,2.56750)(497.48,2.68500)(500.76,2.80250)(504.03,2.92000)(507.31,3.03750)(510.59,3.21500)(513.87,3.36250)(517.14,3.54000)(520.42,3.71750)(523.7,3.95250)(526.97,4.21750)(530.25,4.48500)(533.53,4.83750)(536.81,5.31000)(540.08,5.87000)(543.36,6.49000)(546.64,7.02000)(549.92,7.52250)(553.19,7.96500)(556.47,8.34750)(559.75,8.67250)(563.03,8.96750)(566.3,9.23250)(569.58,9.44000)(572.86,9.67500)(576.13,9.82250)(579.41,9.94000)(582.69,10.06000)(585.97,10.17750)(589.24,10.23500)(592.52,10.35500)(595.8,10.41250)(599.08,10.53000)(602.35,10.59000)(605.63,10.67750)(608.91,10.73750)(612.18,10.79750)(615.46,10.85500)(618.74,10.91500)(622.02,10.97250)(625.29,11.03250)(628.57,11.06250)(631.85,11.09250)(635.13,11.15000)(638.4,11.18000)(641.68,11.24000)(644.96,11.29750)(648.24,11.38750)(651.51,11.50500)(654.79,11.65250)(658.07,11.85750)(661.34,12.06500)(664.62,12.39000)(667.9,12.71500)(671.18,13.12750)(674.45,13.51000)(677.73,13.95250)(681.01,14.36500)(684.29,14.72000)(687.56,15.10250)(690.84,15.39750)(694.12,15.72250)(697.39,15.95750)(700.67,16.22500)(703.95,16.43000)(707.23,16.63750)(710.5,16.81500)(713.78,16.96250)(717.06,17.11000)(720.34,17.25750)(723.61,17.37500)(726.89,17.49250)(730.17,17.61000)(733.45,17.67000)(736.72,17.78750)(740.0,17.84750)(743.28,17.93500)(746.55,18.02250)(749.83,18.08250)(753.11,18.14250)(756.39,18.23000)(759.66,18.29000)(762.94,18.37750)(766.22,18.43750)(769.5,18.49500)(772.77,18.58500)(776.05,18.67250)(779.33,18.73250)(782.61,18.85000)(785.88,18.93750)(789.16,19.02750)(792.44,19.14500)(795.71,19.26250)(798.99,19.38000)(802.27,19.49750)(805.55,19.61750)(808.82,19.73500)(812.1,19.82250)(815.38,19.88250)(818.66,19.94000)(821.93,20.00000)(825.21,20.03000)(828.49,20.03000)(831.76,20.06000)(835.04,20.03000)(838.32,20.03000)(841.6,20.00000)(844.87,19.97000)(848.15,19.91250)(851.43,19.85250)(854.71,19.79250)(857.98,19.67500)(861.26,19.55750)(864.54,19.35000)(867.82,18.73250)(871.09,17.58000)(874.37,16.22500)(877.65,15.69250)(880.92,15.31000)(884.2,14.95500)(887.48,14.57250)(890.76,14.07000)(894.03,13.24500)(897.31,11.94750)(900.59,10.35500)(903.87,8.96750)(907.14,8.02250)(910.42,7.55250)(913.7,7.25750)(916.97,6.99000)(920.25,6.78500)(923.53,6.54750)(926.81,6.37250)(930.08,6.13500)(933.36,5.90000)(936.64,4.39500)(939.92,4.63250)(943.19,3.12750)(946.47,2.92000)(949.75,2.80250)(953.03,2.68500)(956.3,2.56750)(959.58,2.50750)(962.86,2.39000)(966.13,2.33000)(969.41,2.21250)(972.69,2.15250)(975.97,2.06500)(979.24,2.00500)(982.52,1.91750)(985.8,1.85750)(989.08,1.77000)(992.35,1.68250)(995.63,1.62250)(998.91,1.56250)(1002.18,1.47500)(1005.46,1.38750)(1008.74,1.32750)(1012.02,1.24000)(1015.29,1.18000)(1018.57,1.09250)(1021.85,1.03250)(1025.13,0.94500)};\node[below]at(400,0){$ 400 $};\node[below]at(500,0){$ 500 $};\node[below]at(600,0){$ 600 $};\node[below]at(700,0){$ 700 $};\node[below]at(800,0){$ 800 $};\node[below]at(900,0){$ 900 $};\node[below]at(1000,0){$ 1000 $};\node[left]at(360,5.0){$ 20 $};\node[left]at(360,10.0){$ 40 $};\node[left]at(360,15.0){$ 60 $};\node[left]at(360,20.0){$ 80 $};\node[above]at(360,27){$I,\,\text{отн. ед.}$};\node[above,align=center,text width=5cm]at(760,23){Спектр излучения свечи};\node[below]at(750,-4){Длина волны, нм};\end{scope}\end{tikzpicture}\end{center}
}{
\smallskip\noindent\textbf{Решение}\par
При трении эбонитовой палочки о сукно, происходит электризация, в результате которой эбонитовая палочка заряжается отрицательно. При касании палочкой металлического шарика происходит электризация соприкосновением, в результате чего шар заряжается отрицательно.
Исходя из графиков энергии и лампы накаливания, и свечи недостаточно для того, чтобы происходил фотоэффект, то есть вырывание электронов с поверхности цинкового шара, потому что длина волны излучаемого света данными источниками больше максимальной длины волны $\lambda_{\,\text{кр}}.$ Поэтому при освещении лампой накаливания заряд электрометра не изменяется.
Температура воздуха вблизи лампы накаливания меньше, чем вблизи горящей свечи. При нагревании воздуха вблизи шара увеличивается скорость движения молекул. При высокой температуре происходит ионизация воздуха  — появляются положительные и отрицательные ионы. Так как электрометр обладает отрицательным зарядом, то в результате притяжения между отрицательно заряженным шариком и положительными ионами происходит уменьшение заряда электрометра.
}

\end{document}
